% =======================================================================
% HJM MONTE CARLO IMPLEMENTATION (HAZARD RATES)
% =======================================================================

\section{Monte Carlo Implementation of the HJM Hazard Model}

\paragraph{Note on Sources.}
The Monte Carlo formulation in this section follows the structure of 
\cite{glasserman2004} \textit{Monte Carlo Methods in Financial Engineering} 
, Chapter~3, and in particular Sections~3.6 and~3.6.1.  
The original presentation concerns interest-rate term-structure models, 
and has been translated here into the credit-risk setting.  
The notation has been adjusted, but the foundations are based on Glasserman's work.

The previous subsection introduced the continuous time HJM framework for
forward hazard rates $\gamma(t,T)$ and showed how it preserves today’s survival
curve $S(0,T)$ in expectation. We now translate that theory into a concrete
Monte Carlo recipe that can be coded directly.

The guiding principles are the same as in the classical HJM interest rate
literature, but applied to credit:

\begin{itemize}
    \item Work with the \emph{forward hazard surface} $\gamma(t,T)$, not directly
          with $h(t)$.
    \item Initialise a \emph{discrete} forward hazard surface so that it matches
          the deterministic survival curve exactly on the simulation grid.
    \item Evolve the surface in time using a one factor volatility structure and
          a discrete drift chosen to keep survival probabilities arbitrage free
          on the grid.
    \item Read the diagonal to obtain the instantaneous hazard $h(t)$ along each
          path, then integrate $h(t)$ to obtain survival factors and default
          times.
\end{itemize}

Throughout this section we keep the short rate deterministic, so all discount
factors $D(0,t)$ can be precomputed and taken outside the simulation. The Monte
Carlo part of the model carries only \emph{credit} randomness. 

\subsection{Time and Maturity Grids}

We use a single time grid for both simulation time and credit maturities:
\[
0 = t_0 < t_1 < \dots < t_N = T_{\max},
\qquad
\Delta t_i = t_{i+1} - t_i.
\]

In practice:

\begin{itemize}
    \item Choose $T_{\max}$ to be at least the final maturity of all instruments
          in the portfolio.
    \item Use a relatively fine grid around coupon dates and potential call or
          knock out dates, and a coarser grid in between.
    \item For simplicity we let both arguments of $\gamma$ live on this same
          grid, so we will approximate $\gamma(t_i, T_j)$ only for $0 \le i \le j
          \le N$.
\end{itemize}

We denote the discrete approximation by $\widehat{\gamma}(t_i, t_j)$, or more
compactly by $\gamma_{i,j}$ when the meaning is clear.

\subsection{Precomputation of the Deterministic Survival and Discount Curves}

From the parametric curves $r(t)$ and $h(t)$ we compute

\[
D(0,t)
  = \exp\!\left(-\int_0^t r(u)\,du\right), 
\qquad
S(0,t)
  = \exp\!\left(-\int_0^t h(u)\,du\right),
\]

for all grid points $t = t_0,\dots,t_N$. These integrals are closed form for the
exponential specifications used in this project, so this step is analytic and
cheap.

\begin{itemize}
    \item The discount factors $D(0,t_i)$ will be used only for present valuing
          cashflows and do not interact with the Monte Carlo evolution.
    \item The survival probabilities $S(0,t_i)$ will be used to initialise the
          forward hazard surface in a way that is exactly consistent with the
          deterministic curve.
\end{itemize}

\subsection{Discrete Initial Forward Hazard Surface}

In continuous time the relationship between the forward hazard curve and the
survival function at time $t$ is
\[
S(t,T)
  = \exp\!\left(-\int_t^T \gamma(t,u)\,du\right).
\]

On the discrete grid we approximate this by
\[
\widehat{S}(t_i, t_j)
  = \exp\!\left(-\sum_{l=i}^{j-1} \gamma_{i,l}\,\Delta t_l\right),
  \qquad j > i,
\]
with the convention $\widehat{S}(t_i, t_i) = 1$.

To ensure that the simulated model exactly matches the deterministic survival
curve at $t = 0$ we demand
\[
\widehat{S}(0, t_j) = S(0, t_j),
\qquad j = 0,\dots,N.
\]
This leads to the recursion
\[
\widehat{S}(0, t_{j+1})
 = \widehat{S}(0, t_j)
   \exp\!\left(-\gamma_{0,j}\,\Delta t_j\right),
\]
hence
\[
\gamma_{0,j}
  = -\frac{1}{\Delta t_j}
     \log\frac{S(0, t_{j+1})}{S(0, t_j)},
  \qquad j = 0,\dots,N-1.
\]

\paragraph{Interpretation}
The discrete forward hazard $\gamma_{0,j}$ is the \emph{average} of the
continuous forward hazard over $[t_j,t_{j+1}]$. Directly sampling
$\gamma(0,t_j)$ from the continuous formula would not preserve the survival
curve on the grid. Matching the log ratio
$S(0,t_{j+1})/S(0,t_j)$ is therefore essential.

\subsection{Discretising the Volatility Surface}

The continuous model specifies a one factor exponential volatility structure
\[
\sigma_h(t,T) = \sigma_h\, e^{-c_h(T - t)}.
\]
On the discrete grid we simply sample this at the grid points:
\[
\widehat{\sigma}_h(t_i, t_j)
  = \sigma_h\, e^{-c_h(t_j - t_i)},
  \qquad j \ge i.
\]

For simulation it is convenient to write
\[
\sigma_{i,j} := \widehat{\sigma}_h(t_i, t_j),
\]
and to use a single scalar Brownian motion per path.

\subsection{Discretising HJM Drift on the Grid}

In continuous time the HJM drift restriction ensures that survival curves
remain arbitrage free. On the discrete grid we want an analogous property:
for each maturity $T_j$ the survival value $S(t_i, t_j)$ should evolve as a
martingale when discounted appropriately.

The standard HJM trick is to construct the discrete drift from the volatility
surface directly. In style of \cite{glasserman2004} (Specifically 3.6.2, The discrete drift), define for fixed $i$ the cumulative
volatility weights
\[
w_n
  = \sum_{l=i}^n \sigma_{i-1,l}\,\Delta t_l,
  \qquad n = i,\dots,N-1,
\]
with $w_{i-1} := 0$.

The arbitrage free discrete drift at time step $t_{i-1} \to t_i$ is then given by
\[
\widehat{\alpha}_h(t_{i-1}, t_j)
  = \mu_{i-1,j}
  = \frac{w_j^2 - w_{i-1}^2}{2\,\Delta t_j},
  \qquad j = i,\dots,N-1.
\]

This is the exact discrete counterpart of the continuous time HJM drift
condition. It ensures that the discrete survival curve implied by
$\gamma_{i,\cdot}$ is consistent with no arbitrage on the grid.

\paragraph{Important warning.}
A naive Euler discretisation of
\[
d\gamma(t,T) = \alpha_h(t,T)\,dt + \sigma_h(t,T)\,dW_t
\]
at each maturity $T$ separately breaks the discrete martingale property and
causes slow drift in the survival curve as $i$ increases. The construction
above explicitly couples all maturities through the cumulative weights $w_n$
and avoids this problem.

\subsection{Forward Hazard Update per Time Step}

Given $\gamma_{i-1,j}$ for $j \ge i-1$ at time $t_{i-1}$ we update the forward
hazards to $t_i$ as follows.

\begin{enumerate}
    \item Compute the volatility slice
          \[
          \sigma_{i-1,j} = \sigma_h\, e^{-c_h(t_j - t_{i-1})},
          \qquad j = i,\dots,N-1.
          \]
    \item Form the cumulative volatility weights
          \[
          w_n
            = \sum_{l=i}^n \sigma_{i-1,l}\,\Delta t_l,
            \qquad n = i,\dots,N-1,
          \]
          and set $w_{i-1} = 0$.
    \item Compute the discrete drift
          \[
          \mu_{i-1,j}
            = \frac{w_j^2 - w_{j-1}^2}{2\,\Delta t_j},
            \qquad j = i,\dots,N-1.
          \]
    \item Draw a single standard normal $Z_i \sim \mathcal{N}(0,1)$ for this
          time step.
    \item Update the forward hazards
          \[
          \gamma_{i,j}
            = \gamma_{i-1,j}
              + \mu_{i-1,j}\,\Delta t_{i-1}
              + \sigma_{i-1,j}\,\sqrt{\Delta t_{i-1}}\,Z_i,
            \qquad j = i,\dots,N-1.
          \]
\end{enumerate}

Values with $j < i$ correspond to maturities that have already passed, so they
are no longer part of the forward surface.

% =======================================================================
% 3.7  Pathwise Survival Under HJM Hazard Dynamics
% =======================================================================

\subsection{Pathwise Survival Under HJM Hazard Dynamics}

In many reduced-form credit models, default times are simulated explicitly by
drawing an exponential random variable and comparing it with the integrated
hazard process. While this method is theoretically correct, it produces
\emph{discontinuous} payoff behaviour across paths: cashflows abruptly drop to
zero once a simulated default time, $\tau$, falls before maturity. For instruments
with embedded optionality (e.g.  callable coupon structures, range-accrual
credit-linked notes, and digital redemption features) this “hard default” logic
leads to numerical kinks and unstable Monte Carlo convergence.

In this project we instead adopt a \emph{pathwise survival weighting} approach.
Rather than simulating default times, each Monte Carlo path produces a
\emph{forward hazard surface}
\[
\gamma_{i,j} := \widehat{\gamma}(t_i, t_j),
\qquad j \ge i,
\]
from which a full pathwise survival curve is constructed.

Given the forward hazards at simulation time $t_i$, the model-implied survival
from $t_i$ to $t_j$ is
\[
S(t_i, t_j)
    = \exp\!\left(
        - \sum_{\ell=i}^{j-1}
           \gamma_{i,l}\,\Delta t_\ell
      \right).
\]

This expression is evaluated \emph{on every path}. Consequently:

\begin{itemize}
    \item No default time is drawn.
    \item No cashflow is ever forcibly set to zero.
    \item All valuation is performed under smooth, differentiable
          survival weights $S(t_i,t_j)\in(0,1]$.
\end{itemize}

For each Monte Carlo path $p$, the contribution of a cashflow $c_j$ payable at
$t_j$ is
\[
\mathrm{PV}^{(p)}(c_j)
    = D(0,t_j)\,c_j\,S^{(p)}(t_i, t_j).
\]
Averaging over paths yields an unbiased Monte Carlo estimator of the true
reduced-form value. As we will see later, this smooth
treatment of survival probability is essential for structured CLN features,
ensuring stable simulation behaviour and eliminating numerical discontinuities.


% =======================================================================
% 3.8  From Forward Hazards to Instantaneous Hazard and Model CDS Levels
% =======================================================================

\subsection{From Forward Hazards to Instantaneous Hazard and Model CDS Levels}

As described in the previous chapter, the HJM hazard model evolves the \emph{forward hazard surface}
$\gamma(t,T)$ across simulation time. The instantaneous hazard rate at
time $t_i$ is obtained from the diagonal of this surface,
\[
h_i = \gamma_{i,i}.
\]
This instantaneous intensity determines the short-term default risk on
each Monte Carlo path.

For exotic credit-linked notes, coupons or
embedded options are often linked to the prevailing level of credit
risk, typically expressed in terms of a CDS spread. In a full
market-accurate implementation, the par CDS spread would be computed
using protection and premium leg present values, accounting for
discounting, accrual on default, payment conventions, and the term
structure of survival probabilities.

Such detail is unnecessary for the numerical experiments in this
tutorial. Instead, the \emph{model CDS spread} at time $t_i$ is defined
pathwise by
\[
s_i := (1 - R)\,h_i = (1 - R)\,\gamma_{i,i}.
\]

This quantity serves as the credit-risk state variable inside CLN
payoffs. All CDS-linked features are evaluated using
$s_i$. Since $h_i$ depends smoothly on the evolved forward hazards, the
resulting model CDS level reacts smoothly to volatility and curve
shifts, ensuring stable Monte Carlo behaviour.

\subsubsection*{Remark on Modelling Choice}

This simplified approach is appropriate here because the focus of the
tutorial is on the Monte Carlo implementation of the HJM hazard model,
rather than on exact replication of market CDS conventions. Should a
more realistic CDS calibration be required, the same forward-hazard
paths may be used to compute par spreads using full survival-based
leg valuations. That extension is conceptually straightforward but adds
substantial notational overhead without affecting the numerical
principles developed in this section.


\subsection{Exotic Range--Accrual CLN Payoff}

This section introduces the exotic payoff used throughout Chapter~4 to illustrate
variance--reduction methods. The product is a range--accrual credit--linked note
in which the final redemption amount is scaled by the proportion of
coupon--observation dates on which the model--implied CDS spread stays inside a
fixed corridor. The structure preserves the standard coupon, recovery, and
redemption legs of a risky bond while introducing controlled path dependence.

Let $t_1,\dots,t_N$ denote the coupon observation dates. Along each simulated
hazard--rate path, the HJM model produces a model--implied CDS spread $s(t_i)$.
Fix a corridor $[L,U]$ and define
\[
I_i = 1\{L \le s(t_i) \le U\},
\qquad
A = \frac{1}{N}\sum_{i=1}^N I_i.
\]
The factor $A$ measures the fraction of times the spread remains within the
corridor. At maturity, the notional is scaled by this factor, producing a smooth
and intuitive form of path dependency.

Let $D(0,t)$ and $S(0,t)$ denote the pathwise discount factor and pathwise
survival probability along the simulated hazard--rate path.

\subsubsection*{Coupon Leg}

The coupon leg follows the standard reduced--form representation:
\[
\text{CouponLeg}
= \sum_{i=1}^N c\, D(0,t_i)\, S(0,t_i),
\]
where $c$ is the contractual coupon. No exotic modification is applied to this
component.

\subsubsection*{Recovery Leg}

In reduced--form credit models, default can occur at any time, and its timing is
governed by the hazard rate $h(t)$. The key quantity is the \emph{default
density},
\[
h(t)\, S(0,t)\, dt,
\]
which represents the probability of default in the small interval around $t$.
Multiplying this by the recovery amount $R$ and the discount factor $D(0,t)$
gives the time--$t$ contribution of recovery to the overall price. In continuous
time, the recovery leg is
\[
\text{RecoveryLeg}
= R \int_0^T D(0,t)\, h(t) S(0,t)\, dt.
\]

For implementation, this integral is approximated on the same grid used for
coupon dates and survival simulation. This maintains conceptual clarity, avoids
the need for finer quadrature, and is adequate for the pedagogical and
variance--reduction focus of this project.

A convenient discrete approximation, which avoids evaluating $h(t)$ directly,
uses the identity
\[
h(u_j)\, S(0,u_j)\, \Delta u_j
\approx S(0,u_{j-1}) - S(0,u_j),
\]
expressing default over the interval $[u_{j-1},u_j]$ purely in terms of survival
probabilities. On the grid $0 = u_0 < u_1 < \dots < u_N = T$, the recovery leg
is therefore
\[
\text{RecoveryLeg}
= R \sum_{j=1}^{N} D(0,u_j)\, \bigl(S(0,u_{j-1}) - S(0,u_j)\bigr).
\]

\subsubsection*{Redemption Leg}

Only the redemption amount is affected by the range--accrual feature. Instead of
paying the full notional at maturity, the product pays
\[
\text{RedemptionLeg}
= A \cdot \mathrm{Notional} \cdot D(0,T)\, S(0,T).
\]
If the CDS spread is always inside the corridor, $A = 1$ and the standard CLN
redemption is recovered; if it spends significant time outside the corridor, the
redemption amount is proportionally reduced.

\subsubsection*{Full Exotic Payoff}

Combining all components, the payoff for a single Monte Carlo path is
\[
\text{Payoff}
= \text{CouponLeg}
+ \text{RecoveryLeg}
+ \text{RedemptionLeg}.
\]

This formulation maintains the classical risky--bond decomposition while
introducing a smooth and interpretable path--dependent element. Because the
exotic modification interacts closely with the vanilla CLN structure but varies
nonlinearly with spread dynamics, it serves as an ideal payoff for demonstrating
the effect of common random numbers, control variates, and antithetic variates.

This completes the numerical cookbook for the HJM hazard model. The next section will discuss variance reduction 
techniques in detail. 